\documentclass[11pt]{article}
\usepackage[margin=1in]{geometry}
\usepackage{hyperref}
\usepackage{booktabs}
\usepackage{xcolor}
\usepackage{enumitem}
\hypersetup{colorlinks=true, urlcolor=blue, linkcolor=black}

\title{JOLT: Literature Justification for a Legged-Motion\\
Robustness Benchmark in Outdoor Scene Segmentation}
\author{JOLT project}
\date{June 2026}

\begin{document}
\maketitle

\begin{abstract}
This document answers one question before any downstream code is written:
\emph{is there sufficient literature to justify the use case behind JOLT?} The
verdict is \textbf{yes for justification, with an important caveat for novelty}.
The corruption-robustness methodology is well established and directly transferable,
and the off-road/legged specialization is genuinely under-served --- but several
adjacent ideas (corruption-robustness \emph{for segmentation}, motion-blur-robust
segmentation, IMU-conditioned deblurring) are \emph{already owned} in the on-road /
image-restoration literature. The defensible contribution is therefore a
\emph{domain-and-metric specialization}, not a methodological first. We document
the precedents, the genuine whitespace, the threats to novelty, and the evidence
that makes the empirical gate (\S\ref{sec:gate}) the deciding experiment.
\end{abstract}

\section{The claim under test}
JOLT rests on one empirical and one positioning claim:
\begin{enumerate}[noitemsep]
  \item \textbf{Empirical:} legged-motion corruption (motion blur, rolling shutter,
        viewpoint jitter) materially degrades SOTA terrain segmentation. \emph{(Settled
        by the gate, not by literature --- \S\ref{sec:gate}.)}
  \item \textbf{Positioning:} no existing benchmark or method \emph{owns} this for
        off-road / legged terrain. \emph{(Assessed here.)}
\end{enumerate}

\section{Precedent: the methodology is established (this is a strength)}
\label{sec:precedent}
The relative-robustness framing is mature and transferable, which de-risks the method:
\begin{itemize}[noitemsep]
  \item \textbf{ImageNet-C} \cite{imagenetc} introduced parameterized corruptions at
        graded severities and the \emph{relative} robustness metric. This is the exact
        template JOLT's retention $R = m(c,s)/m_0$ follows, and the reason coarse
        ground truth cancels.
  \item \textbf{Robustness of semantic segmentation to common corruptions}
        \cite{kamann} is the direct segmentation analogue (Cityscapes / PASCAL VOC /
        ADE20K), and it covers blur corruptions: the paper names PSF blur, and whether it
        includes \emph{motion} blur could not be confirmed from the abstract or the
        available snippets. This both validates the approach and is the strongest novelty
        threat (\S\ref{sec:threats}).
  \item \textbf{Motion-blur-robust segmentation via class-centric augmentation}
        \cite{classcentric} (CVPR 2023) shows motion-blur robustness for segmentation
        is itself an active, partly-owned topic --- but on-road.
  \item \textbf{LiDAR segmentation robustness benchmarks} \cite{robo3d} and the
        \textbf{ACDC} adverse-condition driving benchmark \cite{acdc} extend the
        robustness-benchmark idea to other modalities / real conditions --- again,
        \emph{urban/driving}, not off-road and not gait-matched.
\end{itemize}

\section{The genuine whitespace (this is the contribution)}
\label{sec:whitespace}
Cross-referencing the corruption-robustness literature (\S\ref{sec:precedent}, all
urban/driving) against the off-road perception literature (which exists but is
accuracy- or generalization-focused, not corruption-focused) leaves a clear gap:
\begin{itemize}[noitemsep]
  \item \textbf{No corruption-robustness benchmark targets off-road / unstructured
        terrain.} Off-road work \cite{offseg,ganav,crossdataset} optimizes clean
        accuracy or cross-dataset generalization; robustness benchmarks
        \cite{kamann,acdc,robo3d} target cities/roads. The intersection is empty.
  \item \textbf{No benchmark ties corruption severity to legged gait.} Existing
        severity ladders are generic; none are calibrated to gait frequency
        $\times$ exposure or to legged ego-motion statistics.
  \item \textbf{The noisy-GT argument is independently supported.} Off-road labels
        are documented as coarse/ambiguous (e.g. coarse-label off-road work
        \cite{coarse}), which is precisely why an accuracy target is weak and a
        \emph{retention} target (where label noise cancels) is the right reframe.
  \item \textbf{Legged ground-level perception is rising but corruption-blind.}
        Recent legged datasets/benchmarks \cite{grandtour,oneocc,robosense} establish
        the platform's importance. Legged gait blur has been studied for feature
        detection \cite{gultekin} and for panoramic occupancy prediction
        \cite{panommocc,oneocc}, yet none of these stress-test segmentation under
        synthesized legged-motion corruption. This narrows the whitespace compared with
        earlier drafts.
\end{itemize}
\noindent\textbf{Net:} the defensible novelty is the \emph{off-road $\times$
legged-gait-matched $\times$ retention-metric} specialization, plus releasing the
parameterized transforms as a reusable benchmark.

\section{Threats to novelty --- read before committing}
\label{sec:threats}
The literature also \emph{narrows} the claim. The JOLT doc's phrase ``no benchmark or
method owns it'' is too strong as written; the precise, defensible version is ``no
benchmark owns it \emph{for off-road / legged terrain}.'' Specifically:
\begin{itemize}[noitemsep]
  \item Do \textbf{not} claim the first corruption-robustness benchmark for
        segmentation --- \cite{kamann} owns it.
  \item Do \textbf{not} claim the first motion-blur-robust segmentation ---
        \cite{classcentric} owns it (on-road).
  \item Do \textbf{not} claim the first IMU/proprioception-conditioned deblurring ---
        gyroscope- and IMU-aided deblurring already exist \cite{mustaniemi,gyrodeblur,imumesh},
        and explicitly mention legged robots. The novelty would be \emph{coupling it
        to terrain segmentation robustness}, not the deblurring idea itself.
\end{itemize}
Mitigation: frame as a \emph{domain benchmark + targeted method}, lead with the
gait-matched calibration and the off-road retention story, and (decisively) show that
even the foundation-model ceiling degrades.

\section{Why the gate is still the deciding experiment}
\label{sec:gate}
Literature can justify the \emph{framing} but cannot settle the \emph{empirical}
claim: foundation models (DINOv2, SAM2) trained at scale may already be robust to
these corruptions, in which case the benchmark has no signal and the thread should
pivot to cross-dataset generalization. Two outcomes maximize the contribution and are
exactly what the gate measures:
\begin{enumerate}[noitemsep]
  \item a \textbf{large, monotone} retention drop with severity, and
  \item the drop persisting for \textbf{SAM2 (the ceiling)}, not just the small
        frozen-DINOv2 head --- which rules out ``it's just a weak head.''
\end{enumerate}
The gate harness (\texttt{jolt/scripts/gate.py}) computes retention and boundary-F
retention per corruption $\times$ severity and prints an advisory verdict.

\section{Verdict}
\textbf{Sufficient literature to justify the use case: yes.} The methodology is
proven and transferable; the off-road/legged specialization is a real, unoccupied
niche; and the noisy-label premise is independently supported. \textbf{But} the
novelty must be scoped as a domain-and-metric specialization, because the generic
versions of the benchmark and the candidate method are already published. Venue
implication: this supports a Datasets-\&-Benchmarks / robotics-venue submission
strongly, and a CV main-track submission only if the gate passes \emph{and} a
targeted method beats baselines on retention across multiple off-road datasets.

\begin{thebibliography}{99}
\bibitem{imagenetc} D. Hendrycks and T. Dietterich.
  \emph{Benchmarking Neural Network Robustness to Common Corruptions and Perturbations.}
  ICLR 2019. \url{https://arxiv.org/abs/1903.12261}
\bibitem{kamann} C. Kamann and C. Rother.
  \emph{Benchmarking the Robustness of Semantic Segmentation Models with Respect to
  Common Corruptions.} IJCV 129 (2021), 462--483.
  \url{https://link.springer.com/article/10.1007/s11263-020-01383-2}
\bibitem{classcentric} Aakanksha and A. N. Rajagopalan.
  \emph{Improving Robustness of Semantic Segmentation to
  Motion-Blur Using Class-Centric Augmentation.} CVPR 2023.
  \url{https://cvpr.thecvf.com/virtual/2023/poster/21045}
\bibitem{robo3d} L. Kong et al. \emph{Benchmarking the Robustness of LiDAR Semantic
  Segmentation Models (Robo3D).} \url{https://arxiv.org/pdf/2301.00970}
\bibitem{acdc} C. Sakaridis et al. \emph{ACDC: The Adverse Conditions Dataset with
  Correspondences for Semantic Driving Scene Understanding.} ICCV 2021.
  \url{https://acdc.vision.ee.ethz.ch/}
\bibitem{gyrodeblur} \emph{Gyroscope-Assisted Motion Deblurring Network.} 2024.
  \url{https://arxiv.org/abs/2402.06854}
\bibitem{imumesh} \emph{IMU-aided Adaptive Mesh-Grid Based Video Motion Deblurring.}
  PeerJ CS, 2024. \url{https://peerj.com/articles/cs-2540/}
\bibitem{coarse} \emph{COARSE: Collaborative Pseudo-Labeling with Coarse Real Labels
  for Off-Road Semantic Segmentation.} 2025. \url{https://arxiv.org/html/2503.03947v1}
\bibitem{crossdataset} \emph{Cross-Dataset Evaluation of Visual Semantic Segmentation
  Models for Off-Road Autonomous Driving.} 2026.
  \url{https://www.sciencedirect.com/science/article/pii/S0957417426015691}
\bibitem{offseg} K. Viswanath et al. \emph{OFFSEG: A Semantic Segmentation Framework
  for Off-Road Driving.} \url{https://arxiv.org/pdf/2103.12417}
\bibitem{ganav} T. Guan et al. \emph{GANav: Efficient Terrain Segmentation for Robot
  Navigation in Unstructured Outdoor Environments.} IEEE RA-L, 2022.
  \url{https://arxiv.org/abs/2103.04233}
\bibitem{grandtour} \emph{GrandTour: A Legged Robotics Dataset in the Wild for
  Multi-Modal Perception and State Estimation.} 2026.
  \url{https://arxiv.org/pdf/2602.18164}
\bibitem{oneocc} \emph{OneOcc: Semantic Occupancy Prediction for Legged Robots with a
  Single Panoramic Camera.} 2025. \url{https://arxiv.org/pdf/2511.03571}
\bibitem{robosense} \emph{RoboSense: Large-scale Dataset and Benchmark for Egocentric
  Robot Perception and Navigation in Crowded and Unstructured Environments.} 2024.
  \url{https://arxiv.org/abs/2408.15503}
\bibitem{rugd} M. Wigness et al. \emph{A RUGD Dataset for Autonomous Navigation and
  Visual Perception in Unstructured Outdoor Environments.} IROS 2019.
\bibitem{rellis} P. Jiang et al. \emph{RELLIS-3D Dataset: Data, Benchmarks and Analysis.}
  \url{https://arxiv.org/abs/2011.12954}
\bibitem{goose} \emph{GOOSE Dataset.} \url{https://goose-dataset.de/}
\bibitem{dinov2} M. Oquab et al. \emph{DINOv2.} \url{https://arxiv.org/abs/2304.07193}
\bibitem{gultekin} G\"ultekin and Saranl\i. \emph{Feature detection performance based
  benchmarking of motion deblurring methods: applications to vision for legged robots.}
  Image and Vision Computing 82 (2019), 26--38. (Author initials not recorded.)
\bibitem{panommocc} \emph{Panoramic Multimodal Semantic Occupancy Prediction for
  Quadruped Robots.} arXiv:2603.13108, 2026. \url{https://arxiv.org/abs/2603.13108}
\bibitem{mustaniemi} Mustaniemi et al. \emph{Gyroscope-Aided Motion Deblurring with Deep
  Networks.} WACV 2019. arXiv:1810.00986. (Author initials not recorded.)
\end{thebibliography}
\end{document}
